%=======================================
% ~800-word SOP body (condensed from Rhea_SOP_body.tex), shared by the
% UChicago Harris (800-word limit) and Northwestern (<1000 words) SOPs.
% Edits to Rhea_SOP_body.tex do NOT flow through here -- update both.
%=======================================
\shead{Research interests} I am applying to pursue a PhD in \program{} at \school, where I intend to study questions at the intersection of public and labour economics. I am specifically focused on the effects of public assistance and social program design on labour market outcomes.

\shead{Academic background} I completed a Bachelor of Commerce (Honours) majoring in Economics and Finance at the University of Melbourne, achieving First Class Honours and receiving numerous awards for my academic performance. My Honours thesis, which evaluated the effects of unconventional monetary policy on financial stability in an emerging market context, was published as a central bank working paper. I also achieved excellent grades in mathematics subjects, including Real Analysis and Linear Algebra.

\shead{Pre-doctoral and professional experience} After my undergraduate degree, I worked as a senior analyst for three years at the Reserve Bank of Australia, Australia's equivalent of the Federal Reserve. I now work as a predoctoral researcher for Associate Professor Manasi Deshpande on topics in labour and public economics. Our current project examines the labour supply effects of the NDIS, a first-of-its-kind disability support program where the federal government provides individuals with disabilities a budget to spend in the private market on services like personal care and therapy. We exploit the staggered rollout of the NDIS to estimate the causal effects on wages and labour supply. We find that the marginal entrants induced into disability work are low-income, lower-skilled women who obtain modest earnings increases from entry, while skilled occupations such as therapists experience higher wages and lower entry. Through this research I have gained significant experience working with large administrative datasets and applying causal inference methods such as staggered rollout event studies, difference-in-differences and shift-share instruments, which I believe have strongly prepared me to conduct independent research.

\shead{Current research} I am interested in understanding the relative merits of different policy approaches to increase labour supply in high-demand sectors. A common mechanism is to raise wages (e.g., \cite{dalbo_finan_rossi2013}), but its effect on employment is ambiguous in equilibrium, since it would also raise firm labour costs. In addition, if there are high fixed costs to reskilling in these areas, the reservation wage required to compensate for this education may be too high. As such, another potential approach could be to subsidise education, thereby reducing fixed costs of entry \citep{schoonover2026}. The effectiveness of either approach depends on the labour supply elasticity of earnings versus entry costs. \\~\\I believe two policy programs in Australia could help answer this question. Firstly, since 2019 the state of Victoria has offered free education in specific courses through `Free TAFE', to incentivise students to pursue careers in priority occupations; as a result it had broad eligibility, with residency as the main condition. Secondly, a national 15 per cent wage increase for aged care workers from 2023 \citep{fwc_agedcare2022} is similarly targeted to a specific sector. Unlike other minimum wage changes, the Australian Federal Government is funding this increase by effectively subsidising firms (\href{https://www.health.gov.au/ministers/the-hon-mark-butler-mp/media/budget-2023-24-delivering-the-largest-ever-pay-rise-to-aged-care-workers?language=en&utm_source=miragenews&utm_medium=miragenews&utm_campaign=news}{Australian Government Department of Health (2023)}), so firms may have less incentive to change their input mix. \\~\\I have commenced this research starting with the education subsidy, identifying the causal effect of Free TAFE on enrolment and characterising the marginal enrollee, which is useful to compare against the marginal entrant under minimum wage increases. Using Australian administrative data and a difference-in-differences comparing Victoria to other states, I find that Free TAFE generated a large and immediate increase in enrolments, particularly among women, working-aged cohorts, employed workers, and for more advanced courses. Constructing a plausible `risk set' of likely enrollees based on ex-ante characteristics, I find that young women previously on income support have the strongest enrolment first stage, which I have verified using a distance-to-college instrument \citep{card1995}. I am now estimating employment and wage responses to Free TAFE, and plan to compare these against entry responses to the aged care minimum wage increase.

\shead{Future research} My work on Free TAFE sits within a broader agenda on how labour and welfare outcomes vary by the different types of support public assistance programs provide, especially beyond income assistance. For instance, another input into the labour supply decision over and above wages and fixed costs of working is family planning (e.g., \cite{gallen2026}). I intend to examine a `free IVF' initiative in Victoria, where IVF is now free at a range of public fertility clinics with no means test. This policy lowers one of the `fixed costs' of trying to have a child later, rather than replacing lost income from childbirth, and so could generate different childbirth and labour market decisions than parental income assistance.
